\subsection{CHANGE OF ORDER OF INTEGRATION}

Let \(I = [a, b]\) and \(A = [c, d]\) be two compact intervals in \(R\) ; let \(f\) be a continuous function on \(I \times A\) with values in a complete normed space \(E\) over \(R\) ; by prop. 2 of II, p.~69, \(\int_{a}^{b} f(x, \alpha) dx\) is a continuous function of \(\alpha\) on \(A\) ; its integral \(\int_{c}^{d} \left( \int_{a}^{b} f(x, \alpha) dx \right) d\alpha\) is also denoted, for simplicity, by \(\int_{c}^{d} d\alpha \int_{a}^{b} f(x, \alpha) dx\) .

PROPOSITION 8. If \(\mathbf{f}\) is continuous on \(\mathrm{I} \times \mathrm{A}\) one has

\[
\int_ {c} ^ {d} d \alpha \int_ {a} ^ {b} \mathbf {f} (x, \alpha) d x = \int_ {a} ^ {b} d x \int_ {c} ^ {d} \mathbf {f} (x, \alpha) d \alpha\tag{15}
\]

(``formula for interchanging the order of integration'').

We shall show that, for all \(y \in A\) , one has

\[
\int_ {c} ^ {y} d \alpha \int_ {a} ^ {b} \mathbf {f} (x, \alpha) d x = \int_ {a} ^ {b} d x \int_ {c} ^ {y} \mathbf {f} (x, \alpha) d \alpha .\tag{16}
\]

Since the two sides of (16) are functions of \(y\) , and equal for \(y = c\) , it will suffice to prove that they are differentiable on \(]c\) , \(d[\) and that their derivatives are equal at every point of this interval. If one puts \(\mathbf{g}(\alpha) = \int_{a}^{b} \mathbf{f}(x, \alpha) dx\) , and \(\mathbf{h}(x, y) = \int_{c}^{y} \mathbf{f}(x, \alpha) dx\) , the relation (16) can be written

\[
\int_ {c} ^ {y} \mathbf {g} (\alpha) d \alpha = \int_ {a} ^ {b} \mathbf {h} (x, y) d x.
\]

Now, the derivative of the first term with respect to \(y\) is \(\mathbf{g}(y)\) , while that of the second is \(\int_{a}^{b} \mathbf{h}_{y}'(x, y) dx\) , by II, p.74, corollary, since \(\mathbf{h}_{y}'(x, y) = \mathbf{f}(x, y)\) is continuous on \(I \times A\) ; the two expressions thus obtained are identical.

Suppose now that A = {[}c, d{]} is a compact interval in R, and I an arbitrary interval in R; let f be a continuous function on I × A, with values in E, such that the integral \(\mathbf{g}(\alpha) = \int_{\mathrm{I}} \mathbf{f}(t, \alpha) dt\) is convergent for all \(\alpha \in A\) ; even if \(\mathbf{g}(\alpha)\) is continuous on A one cannot always interchange the order of integration in the integral \(\int_{c}^{d} d\alpha \int_{I} \mathbf{f}(t, \alpha) dt\) , for the integral \(\int_{I} dt \int_{c}^{d} \mathbf{f}(t, \alpha) d\alpha\) may not exist, or it may be different from the integral \(\int_{c}^{d} d\alpha \int_{I} \mathbf{f}(t, \alpha) dt\) (cf.~II, p.~87, exerc. 7). One has, however, the following result:

PROPOSITION 9. If the function \(\mathbf{f}\) is continuous on \(\mathrm{I} \times \mathrm{A}\) , and if the integral \(\int_{\mathrm{I}} \mathbf{f}(t, \alpha) dt\) is uniformly convergent on \(\mathrm{A}\) , then the integral \(\int_{\mathrm{I}} dt \int_{c}^{d} \mathbf{f}(t, \alpha) d\alpha\) is convergent, and one has

\[
\int_ {c} ^ {d} d \alpha \int_ {\mathrm{I}} \mathbf {f} (t, \alpha) d t = \int_ {\mathrm{I}} d t \int_ {c} ^ {d} \mathbf {f} (t, \alpha) d \alpha .\tag{17}
\]

For every compact interval J contained in I, put \(\mathbf{u}_{\mathrm{J}}(\alpha) = \int_{\mathrm{J}} \mathbf{f}(t, \alpha) dt\) . The hypothesis entails that with respect to the filter of sections \(\Phi\) of the directed set \(\Re(\mathrm{I})\) the continuous function \(u_{J}\) converges uniformly on A to \(\int_{I} f(t, \alpha) dt\) ; thus (II, p.~68, prop. 1), \(\int_{c}^{d} d\alpha \int_{J} f(t, \alpha) dt\) has limit \(\int_{c}^{d} d\alpha \int_{I} f(t, \alpha) dt\) with respect to \(\Phi\) ; but, by prop. 8 (II, p.~77), one has

\[
\int_ {c} ^ {d} d \alpha \int_ {\mathrm{J}} \mathbf {f} (t, \alpha) d t = \int_ {\mathrm{J}} d t \int_ {c} ^ {d} \mathbf {f} (t, \alpha) d \alpha .\tag{18}
\]

The preceding result thus means that the integral \(\int_{I} dt \int_{c}^{d} \mathbf{f}(t, \alpha) d\alpha\) is convergent, and on passing to the limit with respect to \(\Phi\) in the relation (18), one obtains (17).

\setcounter{exercisegroup}{0}
\refstepcounter{exercisegroup}
